\subsection{The full scalar subgroup at level four}\label{sec:divisorial-scalar}
This subsection computes the full scalar fixed field at level four and strengthens the ramification obstruction for the associated principal Cartan quotient.

\begin{proposition}\label{prop:divisorial-scalar-level-four}
Use the notation of Proposition~\ref{prop:divisorial-last-cartan}, and take $n=4$. Thus
\[
 A=p^2+p+1,\qquad B=p^2,\qquad N=p^6,
 \qquad w(x)=A,\quad w(y)=B,
\]
where $x=x_4$, $y=y_4$, and $p$ is odd. Write $X,Y$ for their initial forms and put
\[
 S=(1+p\Z_p)/(1+p^4\Z_p)\subseteq\GL_2(\Z/p^4),
 \qquad t=\operatorname{res}(x^B/y^A).
\]
Then
\begin{equation}\label{eq:divisorial-scalar-graded}
 \operatorname{Frac}(\operatorname{gr}_w F^S)
       =k(X^p,Y^{p^2}),
 \qquad e(F/F^S)=p,
 \qquad \kappa(F^S)=k(t^{p^2}).
\end{equation}
In particular $f(F/F^S)=p^2$, and the residue extension is purely inseparable. For the full principal nonsplit Cartan subgroup
\[
 C=(1+p\Z_p[J])/(1+p^4\Z_p[J])
\]
one has
\begin{equation}\label{eq:divisorial-cartan-scalar-bound}
 f(F/F^C)\ge p^2,\qquad e(F/F^C)\le p^4.
\end{equation}
\end{proposition}

We first give the approximation argument needed to pass between the actual fixed fields. For a valued field $K$, write
\[
 \mathcal G(K)=\operatorname{Frac}(\operatorname{gr}_w K).
\]
The degree-zero subfield of its homogeneous elements is its residue field. Every homogeneous element of $\mathcal G(K)$ is the initial form of an element of $K$: a homogeneous quotient can be represented using homogeneous numerator and denominator and then lifted.

\begin{lemma}\label{lem:divisorial-controlled-lift}
Let $K$ be a complete discretely valued field and let $\sigma$ be a finite-order valuation-preserving automorphism. Put $L=K^{\langle\sigma\rangle}$. Suppose that $\sigma-1$ raises every value by at least $d>0$, its degree-$d$ initial operator is a derivation $D$ of $\mathcal G(K)$, and
\[
 \mathcal G(L)=\ker D.
\]
If $c=w(\sigma z-z)>w(z)+d$, there is $a\in L$ such that
\[
 w(z-a)\ge c-d.
\]
The conclusion also holds with $a=z$ when $z$ is fixed.
\end{lemma}

\begin{proof}
The strict inequality implies $D(\operatorname{in}(z))=0$. Choose $a\in L$ with the same initial form as $z$. If the error $b=z-a$ has value below $c-d$, then
\[
 \sigma b-b=\sigma z-z,
 \qquad w(\sigma b-b)>w(b)+d.
\]
Hence $D(\operatorname{in}(b))=0$, and $\operatorname{in}(b)$ lifts to an element of $L$. Adding that element to $a$ strictly increases the error value. Discreteness makes this a finite procedure before the value reaches $c-d$.
\end{proof}

We shall use two elementary facts about graded fields. If $L\subset K$ is finite, then
\begin{equation}\label{eq:divisorial-graded-inequality}
 [\mathcal G(K):\mathcal G(L)]
      =e(K/L)f(K/L)\le [K:L].
\end{equation}
To see this, choose representatives of the value-group cosets and a basis of the residue-field extension. Their homogeneous products form a basis of the graded field extension. Any nontrivial linear relation among their lifts would give a nontrivial relation among the terms of least value, so those lifts are linearly independent over $L$.

Second, a displacement bound on lifts of homogeneous rational generators extends to the complete field. Products, inverses and sums first give the bound on their rational expressions. Every homogeneous initial form has such a lift. Successively subtracting these lifts approximates any element to arbitrarily high value, so continuity gives the bound on every element. The same argument determines the initial operator. Here residues in $k$ have fixed Teichm\"uller representatives in $W(k)$.

\begin{proof}[Proof of Proposition~\ref{prop:divisorial-scalar-level-four}]
Put
\[
 E=w(p)=(p^4-1)p^6,\qquad Q_j=p^{2j},
 \qquad d_j=(Q_j-1)B,\quad d'_j=(Q_j-1)A
 \quad(1\le j\le3).
\]
Let $\sigma_j$ denote scalar multiplication by $1+p^j$. The coordinate calculation in Proposition~\ref{prop:divisorial-last-cartan}, with the height-two leading coefficient normalized to one, gives
\begin{equation}\label{eq:divisorial-scalar-displacements}
 \operatorname{in}([p^j]x)=X^{Q_j},\qquad
 \operatorname{in}([p^j]y)=Y^{Q_j}.
\end{equation}
The formal group law satisfies $G(T,Z)-T=Z(1+\text{terms in }(T,Z))$. Consequently
\[
 \operatorname{in}(\sigma_jx-x)=X^{Q_j},\qquad
 \operatorname{in}(\sigma_jy-y)=Y^{Q_j}.
\]
These substitutions preserve the weighted filtration and induce the identity on its associated graded ring; the inverse substitutions do likewise. They therefore preserve $w$ and extend to $F$.

The least displacement excess on $F$ is $d_j$, since $B<A$. Its degree-$d_j$ operator is
\begin{equation}\label{eq:divisorial-scalar-first-derivation}
 D_j=Y^{Q_j}\frac{\partial}{\partial Y}
       \quad\text{on }\mathcal G(F)=k(X,Y).
\end{equation}
It is a derivation because the product of two positive displacements has strictly higher degree. The preceding continuation argument extends the generator calculation to the complete field. Its kernel is $k(X,Y^p)$.

For odd $p$, the scalar group $S$ is cyclic of order $p^3$, generated by $\sigma_1$. Its subgroups of orders $p^2$ and $p$ are generated by $\sigma_2$ and $\sigma_3$, respectively: the binomial formula gives
\[
 v_p\bigl((1+p)^{p^r}-1\bigr)=r+1.
\]
Write
\[
 L=F^{\langle\sigma_3\rangle},\qquad
 K=F^{\langle\sigma_2\rangle},\qquad M=F^S.
\]
The level action is faithful, so each step in $F/L/K/M$ has degree $p$. These fixed fields are complete for the restricted valuation. The last scalar calculation in Proposition~\ref{prop:divisorial-last-cartan} gives
\begin{equation}\label{eq:divisorial-scalar-last-input}
 \mathcal G(L)=k(X,Y^p)=\ker D_3.
\end{equation}

We will take ordinary $p$th powers as well as formal-group multiples. If $w(z)=r$ and $w(\sigma z-z)=r+s$, then
\begin{equation}\label{eq:divisorial-ordinary-power}
 \operatorname{in}\bigl(\sigma(z^p)-z^p\bigr)
       =\operatorname{in}(\sigma z-z)^p
       \quad\text{provided }E>(p-1)s.
\end{equation}
Indeed the mixed binomial terms have value at least $E+pr+s$, whereas the pure last term has value $pr+ps$. All applications below satisfy this condition because
\[
 A<p^3,\qquad
 E>(p-1)(p^6-1)A.
\]
For the second inequality, its right side is less than
$(p-1)p^9=p^{10}-p^9$, while $E=p^{10}-p^6$. We also have
\begin{equation}\label{eq:divisorial-weight-order}
 B<A<pB.
\end{equation}

Apply Lemma~\ref{lem:divisorial-controlled-lift} to $\sigma_3$, first with $x$ and then with the ordinary power $y^p$. It gives $a,b\in L$ with
\begin{align}
 \operatorname{in}(a)&=X,&
 w(a-x)&\ge C_x:=p^6A-(p^6-1)B,\label{eq:divisorial-first-lift-x}\\
 \operatorname{in}(b)&=Y^p,&
 w(b-y^p)&\ge C_y:=p^7B-(p^6-1)B.\label{eq:divisorial-first-lift-y}
\end{align}
The error bounds exceed $A$ and $pB$, respectively. For $j=1,2$ one has the exact strict inequalities
\begin{align*}
 C_x+d_j-Q_jA&=(p^6-Q_j)(A-B)>0,\\
 C_y+d_j-pQ_jB&=(p-1)(p^6-Q_j)B>0.
\end{align*}
Thus the errors in these lifts do not affect the leading displacements:
\begin{equation}\label{eq:divisorial-middle-displacements}
 \operatorname{in}(\sigma_ja-a)=X^{Q_j},\qquad
 \operatorname{in}(\sigma_jb-b)=Y^{pQ_j}
 \quad(j=1,2).
\end{equation}
In the independent coordinates $X,Z=Y^p$ of $\mathcal G(L)$, inequality~\eqref{eq:divisorial-weight-order} makes the least displacement excess $d'_j$. Its initial derivation is
\[
 D'_j=X^{Q_j}\frac{\partial}{\partial X},\qquad
 \ker D'_j=k(X^p,Y^p).
\]
The continuation argument gives the same bound on all of $L$. Every initial form from $K$ lies in $\ker D'_2$. This kernel has degree $p$ in $\mathcal G(L)$, while $[L:K]=p$. Equation~\eqref{eq:divisorial-graded-inequality} therefore forces
\begin{equation}\label{eq:divisorial-middle-fixed}
 \mathcal G(K)=k(X^p,Y^p).
\end{equation}

There are now controlled representatives for the final scalar action. Apply Lemma~\ref{lem:divisorial-controlled-lift} on $L$, using $\sigma_2$, its shift $d'_2$, and~\eqref{eq:divisorial-middle-fixed}. Starting with $a^p$ and $b$ gives $u,v\in K$ such that
\begin{align}
 \operatorname{in}(u)&=X^p,&
 w(u-a^p)&\ge C_u:=(p^5-p^4+1)A,\label{eq:divisorial-second-lift-x}\\
 \operatorname{in}(v)&=Y^p,&
 w(v-b)&\ge C_v:=p^5B-(p^4-1)A.\label{eq:divisorial-second-lift-y}
\end{align}
Here~\eqref{eq:divisorial-ordinary-power} gives initial displacement $X^{p^5}$ for $a^p$, and~\eqref{eq:divisorial-middle-displacements} gives $Y^{p^5}$ for $b$. The second error exceeds $pB$ because
\[
 C_v-pB=(p^4-1)(pB-A)>0.
\]
To compute $\sigma_1$ on these lifts, use its shift $d'_1$ on $L$. The two error estimates needed are
\begin{align*}
 C_u+d'_1-p^3A&=p^2(p-1)(p^2-1)A>0,\\
 C_v+d'_1-p^3B&=p^2(p^2-1)(pB-A)>0.
\end{align*}
It follows that
\[
 \operatorname{in}(\sigma_1u-u)=X^{p^3},\qquad
 \operatorname{in}(\sigma_1v-v)=Y^{p^3}.
\]
The independent coordinates $U_0=X^p,V_0=Y^p$ of $\mathcal G(K)$ have weights $pA,pB$. The smaller displacement excess is now $p(p^2-1)B$, and the initial derivation is
\[
 D''=V_0^{p^2}\frac{\partial}{\partial V_0},\qquad
 \ker D''=k(U_0,V_0^p)=k(X^p,Y^{p^2}).
\]
The action of $S/\langle\sigma_2\rangle$ on $K$ has order $p$ and fixed field $M$. Its invariant initial forms lie in this kernel, which has degree $p$ in $\mathcal G(K)$. Since $[K:M]=p$, equation~\eqref{eq:divisorial-graded-inequality} forces
\[
 \mathcal G(M)=k(X^p,Y^{p^2}).
\]

The two final homogeneous generators have weights $pA,p^2B$, whose greatest common divisor is $p$. Thus the value group of $M$ is $p\Z$, and $e(F/M)=p$. Since $\gcd(A,pB)=1$, its degree-zero subfield is generated by
\[
 \frac{(X^p)^{pB}}{(Y^{p^2})^A}
      =\left(\frac{X^B}{Y^A}\right)^{p^2}=t^{p^2}.
\]
Therefore $\kappa(M)=k(t^{p^2})$, and the residue extension from $F$ has degree $p^2$, purely inseparable. This proves~\eqref{eq:divisorial-scalar-graded}.

Finally $C$ preserves $w$, has order $p^6$ and contains $S$, by Proposition~\ref{prop:divisorial-last-cartan}. The inclusion $F^C\subseteq F^S$ gives
\[
 f(F/F^C)\ge [k(t):k(t^{p^2})]=p^2.
\]
Since $[F:F^C]=p^6$, the fundamental inequality gives $e(F/F^C)\le p^4$, proving~\eqref{eq:divisorial-cartan-scalar-bound}.
\end{proof}

\begin{corollary}\label{cor:divisorial-1300-bound}
In Proposition~\ref{prop:divisorial-scalar-level-four}, let $p=5$. If $F^C/K$ is a tame extension of ramification index $12$ with unchanged residue field, then
\[
 e_K\ge1300.
\]
Thus this level-four divisorial Cartan construction cannot yield absolute ramification index $260$.
\end{corollary}

\begin{proof}
Ramification indices multiply, so
\[
 e_K=\frac{w(p)}{12\,e(F/F^C)}
     \ge\frac{(5^4-1)5^6}{12\cdot5^4}
     =1300.
\]
\end{proof}

This obstruction concerns the specified valuation and Cartan quotient. The existence of other unequal-character boundary comparisons at index $260$ remains open. The argument does not construct the tame quotient, the required descent datum, or a perfect-residue specialization.
